% Input-ready conventional-proof component. Labels use the prefix ro:.
% Requires mean_field.tex and the bibliography entry Azuma1967.

\section{Exponential rarity of negative endpoint orders}
\label{ro:section}

The product-limit calculation in Corollary~\ref{mf:product-limit} has
a probabilistic consequence for uniformly permuted deterministic row
clouds. It uses standard bounded-difference concentration in addition
to the full mean-field criterion.

\begin{corollary}[Uniform random orders]\label{ro:random-order}
Fix \(t>0\), \(r\ge1\), and deterministic rows
\(v_{1,N},\ldots,v_{N,N}\in\overline B(0,L)\). Suppose
\[
 \nu_N=\frac1N\sum_{j=1}^N\delta_{v_{j,N}}\Longrightarrow\nu,
 \qquad
 \Sigma=\int vv^T\,d\nu(v),\qquad t\lambda_{\max}(\Sigma)>1.
\]
Let \(\pi_N\) be a uniformly distributed permutation of the \(N\)
row labels, and set
\[
 A_N^{\pi_N}=I_N+(t/N)V_N^{\pi_N}(V_N^{\pi_N})^T.
\]
There exist constants \(\eta,c,C>0\) and
\(N_0\) such that, for \(N\ge N_0\),
\begin{equation}\label{ro:exponential}
 \mathbb P\left\{
 \frac{D_1(A_N^{\pi_N})}{N(N-1)\per A_N^{\pi_N}}<\eta
 \right\}\le Ce^{-cN}.
\end{equation}
In particular the fraction of labelled permutations with nonpositive
endpoint derivative is at most \(Ce^{-cN}\).
Under any joint coupling of the uniform permutations across orders,
the displayed normalized derivative is at least \(\eta\) almost
surely for all sufficiently large \(N\). Independence across orders
is not required.
\end{corollary}
\begin{proof}
Set \(\rho=dx\,\nu(dv)\). For this product ordered limit,
Corollary~\ref{mf:product-limit} gives the strictly positive number
\[
 c_* =\min_{h\in K_\rho}\frac{m_h(4-m_h)}{12}>0.
\]
Take \(\eta=c_*/2\). There is a weak neighborhood \(U\) of
\(\rho\) and an order \(N_1\) such that every bounded-row array
of order \(N\ge N_1\) with ordered empirical measure in \(U\)
has normalized derivative at least \(\eta\). Otherwise, choose
successively smaller neighborhoods of \(\rho\) and increasing
orders to obtain empirical measures converging to \(\rho\), each
with normalized derivative below \(c_*/2\). This contradicts the
liminf bound \(c_*\) from Theorem~\ref{mf:criterion} and
Corollary~\ref{mf:product-limit}.

By the definition of the weak topology on the compact row-position
space, choose finitely many real continuous functions
\(f_1,\ldots,f_J\), \(\epsilon>0\), and
\(B=\max(1,\max_j\|f_j\|_\infty)\) such that
\[
 \left\{\mu:
  \left|\int f_j\,d\mu-\int f_j\,d\rho\right|<\epsilon
                 \text{ for }1\le j\le J\right\}\subseteq U.
\]
For a uniform permutation \(\pi\), define
\[
 X_{j,N}=\frac1N\sum_{i=1}^N f_j(i/N,v_{\pi(i),N}).
\]
Its expectation is
\[
 \mathbb EX_{j,N}
 =\frac1{N^2}\sum_{i,a=1}^N f_j(i/N,v_{a,N})
 \longrightarrow\int f_j\,d\rho.
\]
Indeed the deterministic position-grid measures converge to Lebesgue
measure and \(\nu_N\Longrightarrow\nu\), so their product measures
converge weakly. For large enough \(N\), every expectation above is
within \(\epsilon/2\) of its limit.

Swapping two assigned row labels changes \(X_{j,N}\) by at most
\(4B/N\). Reveal the permutation one position at a time. Conditional
on the already exposed labels, completions for any two choices of
the next label can be coupled by exchanging just those two labels.
Thus the conditional expectations for the two choices differ by at
most \(4B/N\), and each increment of the associated Doob martingale
has absolute value at most \(4B/N\).

The bounded-increment martingale inequality of Azuma
\cite{Azuma1967} now gives
\begin{equation}\label{ro:test-concentration}
 \mathbb P\{|X_{j,N}-\mathbb EX_{j,N}|\ge\epsilon/2\}
 \le2\exp\left(-\frac{\epsilon^2N}{128B^2}\right).
\end{equation}
For clarity, this version follows directly by applying conditional
convexity to a centered increment \(Z\) with \(|Z|\le a\):
\(\mathbb E(e^{sZ}\mid\mathcal F)\le\cosh(sa)
\le e^{s^2a^2/2}\). Iteration and Markov's inequality give the
two-sided bound \(2e^{-u^2/(2Na^2)}\); substitute
\(a=4B/N\) and \(u=\epsilon/2\).

The event in \eqref{ro:exponential} forces the ordered empirical
measure outside \(U\). The union bound in
\eqref{ro:test-concentration} proves the assertion with
\(C=2J\), \(c=\epsilon^2/(128B^2)\), and a sufficiently large
\(N_0\). Finally \(\sum_NCe^{-cN}<\infty\), so the first
Borel--Cantelli lemma gives the almost-sure conclusion under any
joint coupling.
\end{proof}

The constants and the sufficient order in this corollary are
existential: the weak-neighborhood argument does not supply explicit
finite rates for the mean-field approximation. When a deterministic
cloud admits an ordered endpoint counterexample, the corollary shows
that such negative endpoint orders occupy an exponentially small
fraction of its labelled permutations. It makes no assertion of
monotonicity over the full \(q\)-interval for the other orders.
